\documentclass[11pt]{article}
\usepackage{preamble}
\renewcommand{\answer}[1]{}

\begin{document}

\ladderheading{AP Statistics \textperiodcentered\ Unit 1 \textperiodcentered\ Topic 1.2 \textperiodcentered\ B: 2026-09-10 \textperiodcentered\ E: 2026-09-11}{Numbers That Aren't Quantities}

\focusbanner{TODAY'S PROBLEM}

A school newspaper surveyed six students for a story about getting to school. Part of the data table is shown.\par\smallskip Two reporters disagree about the numbers in it. \textbf{Maya:} ``Zip code and Rating are both quantitative --- they're numbers, so we can average them.'' \textbf{Jordan:} ``They're both categorical. Numbers can be labels.''\par
\begin{center}
\textbf{Getting-to-school survey (minutes; 1--5 stars)}\par\smallskip
\begin{tabular}{|l|c|c|c|c|}
\hline
\textbf{Name} & \textbf{Lunch} & \textbf{ZIP} & \textbf{Min} & \textbf{Stars} \\ \hline
\textbf{Aiden} & A & 01902 & 12 & 4 \\ \hline
\textbf{Bea} & B & 01905 & 35 & 2 \\ \hline
\textbf{Cruz} & A & 01902 & 8 & 5 \\ \hline
\textbf{Dev} & C & 01904 & 22 & 3 \\ \hline
\textbf{Elise} & B & 01905 & 41 & 1 \\ \hline
\textbf{Fiona} & A & 01901 & 15 & 4 \\ \hline
\end{tabular}
\end{center}
\begin{firsttakebox}{FIRST TAKE --- before the video (5 min)}
Is a zip code a quantitative variable? One sentence: yes or no, and the reason you'd give Maya.
\workspace{0.9in}
\end{firsttakebox}

\begin{callout}{\IconBook\ VARIABLES (keep on the page)}{calloutgreen}
An \textbf{individual} is one row: the person or thing the data describe. A \textbf{variable} is one column: a
characteristic that changes from individual to individual.\\
\textbf{Categorical} --- the values are category names or group labels (even if written with digits).\\
\textbf{Quantitative} --- the values are numbers for a \emph{measured or counted quantity}, with a unit.
Test: does ``twice as big'' or ``the average'' mean anything? If not, it's a label.
\end{callout}

\newpage
\textbf{Finish after the video:}\par\smallskip

\dokbadge{1}~\textbf{(a)} Who are the individuals in this data set, and what are the four variables?\par
\workspace{0.8in}

\dokbadge{2}~\textbf{(b)} Classify each of the four variables as categorical or quantitative. For each quantitative variable, state its unit.\par
\workspace{1.3in}

\dokbadge{3}~\textbf{(c)} Decide who is right about \textbf{Zip code} and who is right about \textbf{Rating}, using the definition of a quantitative variable in each case. Then name one specific calculation a reader would find meaningless if Maya's classification were used for Zip code, and explain why it is meaningless.\par
\workspace{1.9in}

\begin{sentenceframebox}\raggedright\textbf{Frame:} Zip code is \blank[1in] because a quantitative variable must be a \blank[1.4in], and a zip code is a \blank[1in].\end{sentenceframebox}

\begin{sentenceframebox}\raggedright\textbf{Frame:} A \blank[1.2in] of zip codes would be meaningless because \blank[2.4in].\end{sentenceframebox}

\begin{summaryexitbox}{EXIT --- turn this in}
One thing the video changed about my first take: \hrulefill\par
\workspace{0.4in}
\end{summaryexitbox}

\bankitem{Optional \#1 --- not collected}{\dokbadge{2}~For each variable, write \textbf{C} (categorical) or \textbf{Q} (quantitative); for each Q, give the unit: (1) jersey number of a soccer player; (2) number of goals scored this season; (3) shoe size; (4) preferred position (goalkeeper, defender, midfielder, forward); (5) minutes played in the last match.}{}

\end{document}
